%2multibyte Version: 5.50.0.2960 CodePage: 1252

\documentclass[11pt]{article}
%%%%%%%%%%%%%%%%%%%%%%%%%%%%%%%%%%%%%%%%%%%%%%%%%%%%%%%%%%%%%%%%%%%%%%%%%%%%%%%%%%%%%%%%%%%%%%%%%%%%%%%%%%%%%%%%%%%%%%%%%%%%%%%%%%%%%%%%%%%%%%%%%%%%%%%%%%%%%%%%%%%%%%%%%%%%%%%%%%%%%%%%%%%%%%%%%%%%%%%%%%%%%%%%%%%%%%%%%%%%%%%%%%%%%%%%%%%%%%%%%%%%%%%%%%%%
\usepackage{amsmath}

\setcounter{MaxMatrixCols}{10}
%TCIDATA{OutputFilter=LATEX.DLL}
%TCIDATA{Version=5.50.0.2960}
%TCIDATA{Codepage=1252}
%TCIDATA{<META NAME="SaveForMode" CONTENT="1">}
%TCIDATA{BibliographyScheme=Manual}
%TCIDATA{Created=Monday, September 22, 2003 09:40:11}
%TCIDATA{LastRevised=Wednesday, December 17, 2025 19:42:27}
%TCIDATA{<META NAME="GraphicsSave" CONTENT="32">}
%TCIDATA{<META NAME="DocumentShell" CONTENT="General\SW\Blank - Standard LaTeX Article">}
%TCIDATA{Language=American English}
%TCIDATA{CSTFile=LaTeX article (bright).cst}

\newtheorem{theorem}{Theorem}
\newtheorem{acknowledgement}[theorem]{Acknowledgement}
\newtheorem{algorithm}[theorem]{Algorithm}
\newtheorem{axiom}[theorem]{Axiom}
\newtheorem{case}[theorem]{Case}
\newtheorem{claim}[theorem]{Claim}
\newtheorem{conclusion}[theorem]{Conclusion}
\newtheorem{condition}[theorem]{Condition}
\newtheorem{conjecture}[theorem]{Conjecture}
\newtheorem{corollary}[theorem]{Corollary}
\newtheorem{criterion}[theorem]{Criterion}
\newtheorem{definition}[theorem]{Definition}
\newtheorem{example}[theorem]{Example}
\newtheorem{exercise}[theorem]{Exercise}
\newtheorem{lemma}[theorem]{Lemma}
\newtheorem{notation}[theorem]{Notation}
\newtheorem{problem}[theorem]{Problem}
\newtheorem{proposition}[theorem]{Proposition}
\newtheorem{remark}[theorem]{Remark}
\newtheorem{solution}[theorem]{Solution}
\newtheorem{summary}[theorem]{Summary}
\newenvironment{proof}[1][Proof]{\textbf{#1.} }{\ \rule{0.5em}{0.5em}}
\input{tcilatex}
\input{tcilatex}
\setlength{\textwidth}{6.5in}
\setlength{\textheight}{9in}
\setlength{\topmargin}{-.50in}
\setlength{\oddsidemargin}{.0in}
\setlength{\evensidemargin}{.0in}

\begin{document}


\begin{center}
{\LARGE Assignment 8}
\end{center}

\noindent Blankenau\hspace*{\stretch{0.7500}}

\noindent Economics 905, Fall 2025

\noindent \noindent \textit{Instructions.} You may work in groups and may
compare answers prior to submission. Joint submissions with up to 4 names
are allowed and encouraged. All work should be legible and concise. A due
date will be announced in class.

Consider an economy with a single final good which is produced using a
variety of intermediate goods. Specifically, $n_{t}$ intermediate goods ($%
x_{i},$ $i\in \left\{ 1,2..n_{t}\right\} $) combine with labor ($L$) to
produce a final output good $y$ according to 
\begin{equation*}
y_{t}=AL^{1-\alpha }\overset{n_{t}}{\underset{i=1}{\sum }}x_{it}^{\alpha }
\end{equation*}%
where $A>0$ and $0<\alpha <1.$ The final good is the numeraire good and is
produced competitively by a representative firm. Profits for the producer of
the final good $\left( \pi _{ft}\right) $ are 
\begin{equation*}
\pi _{ft}=y_{t}-\overset{n_{t}}{\underset{i=1}{\sum }}p_{it}x_{it}-w_{t}L
\end{equation*}%
where $p_{it}$ is the relative price of intermediate good $i$ and $w_{t}$ is
the wage. Each intermediate good is unique and produced by only one firm.
Thus there is imperfect competition. The only input to the production of an
intermediate good is $\phi >0$ units of the final good. Thus the profit for
the intermediate firm is%
\begin{equation*}
\pi _{i}=p_{i}x_{i}-\phi x_{i}
\end{equation*}%
There is also a government that taxes labor income at rate $\tau $ and
subsidizes inventions at rate $\varsigma $ where $0\leq \tau ,\varsigma \leq
1.$ The government balances its budget in each period. There is free entry
into the business of being an inventor. Anyone can pay $\eta >0$ to invent a
new good. The government pays $\varsigma \eta $ of this so the private cost
is $\left( 1-\varsigma \right) \eta .$ The private cost is financed by
borrowing (issuing bonds). The owner of the new intermediate good type gets
exclusive rights to produce it forever. A representative agent supplies one
unit of time supplied inelastically in the labor market. The agents
inelastically supplies one unit of labor each period. The agent chooses
consumption and bond holdings to bring into the subsequent period to
maximize lifetime utility, $\underset{t=0}{\overset{\infty }{\text{ }\sum }}%
\beta ^{t}\frac{c_{t}^{\sigma }}{\sigma }$ where $0\leq \beta \leq 1,$ $%
\sigma \leq 1$ and $c_{t}$ is consumption$.$ The period budget constraint is 
$c_{t}+b_{t+1}=w_{t}\left( 1-\tau \right) +\left( 1+r_{t}\right) b_{t}\ $%
where $w_{t}$ is wage per unit of time worked, and $b_{t}$ is bond holdings
at the beginning of period $t$ and $r_{t}$ is the interest rate on bonds in
period $t$.

\begin{enumerate}
\item Firms' problems.

\begin{enumerate}
\item Find the final goods producer's derived demand for the intermediate
good $x_{it}$.

\item Find the profit maximizing price of the intermediate good $x_{it}$ set
by its producer.

\item Find the one-period profits for the producer of an intermediate good.
Find the present discounted value of all profits accruing to the producer of
an intermediate good. Assume the interest rate is constant and recall $%
\underset{t=0}{\overset{\infty }{\text{ }\sum }}\frac{1}{\left( 1+r\right)
^{t}}=\frac{1}{r}.$
\end{enumerate}

\item Agent's problem.

\begin{enumerate}
\item Write down \ the Bellman equation for the agent's optimization problem.

\item Find the first order condition and the updated envelope condition.

\item Consider a balanced growth path where $\frac{c_{t+1}}{c_{t}}=\frac{%
y_{t+1}}{y_{t}}=\frac{n_{t+1}}{n_{t}}=\gamma $. Find the optimal growth rate
of consumption as a function of the interest rate.

\item Find the equilibrium interest rate from the bond market clearing
condition.

\item Find the growth rate as a function of exogenous items and policy
variables.
\end{enumerate}

\item The social planner takes $n_{t}$ as a state variable and chooses $%
n_{t+1}$ and $\hat{x}_{t}$ where $\hat{x}_{t}$ is a vector of the $n_{t}$
intermediate goods available in period $t$. It does this to maximize $%
\overset{\infty }{\underset{t=0}{\sum }}\beta ^{t}\ln c_{t}$ where $c_{t}=%
\overset{n_{t}}{\underset{i=1}{\sum }}\left( x_{it}^{\alpha }-x_{it}\right)
-\eta \left( n_{t+1}-n_{t}\right) .$ Here $x_{it}$ is the $i^{th}$ element
of $\hat{x}_{t}.$

\begin{enumerate}
\item Write down the social planner's Bellman equation and find the first
order condition with respect to $x_{it}.$

\item Use the above first order condition and the fact that all values of $x$
will be the same to write the social planner's Bellman equation with only $%
n_{t}$ and $n_{t+1}$ as endogenous items.

\item Find the growth rate that arises as a solution to the social planner's
problem.

\item Find $\varsigma $ such that $x_{it}$ in a competitive equilibrium
coincides with that from the social planner's problem.
\end{enumerate}

\pagebreak \bigskip
\end{enumerate}

\bigskip

\begin{enumerate}
\item \bigskip From notes 
\begin{equation*}
x_{it}=\left( \frac{\alpha AL^{1-\alpha }}{p_{it}}\right) ^{\frac{1}{%
1-\alpha }}
\end{equation*}

\item Profit for the intermediate firm is 
\begin{equation*}
\pi _{it}=p_{it}x_{it}-\phi x_{it}
\end{equation*}%
The firm chooses price to maximize profit. The first order condition is%
\begin{equation*}
\left( p_{it}-\phi \right) \frac{1}{1-\alpha }\left( \frac{\alpha
AL^{1-\alpha }}{p_{it}}\right) ^{\frac{1}{1-\alpha }-1}\frac{\alpha
AL^{1-\alpha }}{p_{it}^{2}}=\left( \frac{\alpha AL^{1-\alpha }}{p_{it}}%
\right) ^{\frac{1}{1-\alpha }}
\end{equation*}%
so 
\begin{equation*}
p_{it}=p_{i}=\frac{\phi }{\alpha }
\end{equation*}

\item Thus%
\begin{equation*}
y_{t}=AL^{1-\alpha }n\left( \frac{\alpha AL^{1-\alpha }}{\frac{\phi }{\alpha 
}}\right) ^{\frac{\alpha }{1-\alpha }}
\end{equation*}%
\begin{equation*}
x_{it}n=n\left( \frac{\alpha AL^{1-\alpha }}{\frac{\phi }{\alpha }}\right) ^{%
\frac{1}{1-\alpha }}
\end{equation*}%
\begin{equation*}
x_{it}=\left( \frac{\alpha AL^{1-\alpha }}{\frac{\phi }{\alpha }}\right) ^{%
\frac{1}{1-\alpha }}
\end{equation*}%
\begin{equation*}
x_{it}=x_{i}=\left( \frac{\alpha ^{2}AL^{1-\alpha }}{\phi }\right) ^{\frac{1%
}{1-\alpha }}
\end{equation*}%
and profits are 
\begin{equation*}
\pi _{it}=\pi _{i}=\left( \frac{\phi }{\alpha }-\phi \right) \left( \frac{%
\alpha ^{2}AL^{1-\alpha }}{\phi }\right) ^{\frac{1}{1-\alpha }}
\end{equation*}

\item 
\begin{eqnarray*}
&&\left( \frac{\phi }{\alpha }-\phi \right) \left( \frac{\alpha
^{2}AL^{1-\alpha }}{\phi }\right) ^{\frac{1}{1-\alpha }}\left( \underset{t=1}%
{\overset{\infty }{\sum }}\frac{1}{\left( 1+r\right) ^{t}}\right) \\
&&\frac{1}{r}\left( \frac{\phi }{\alpha }-\phi \right) \left( \frac{\alpha
^{2}AL^{1-\alpha }}{\phi }\right) ^{\frac{1}{1-\alpha }}
\end{eqnarray*}

\item 
\begin{equation*}
\upsilon \left( b_{t}\right) =\underset{b_{t+1}}{\max }\left[ w_{t}\left(
1-L_{t}\right) \left( 1-\tau \right) +\left( 1+r_{t}\right) b_{t}-b_{t+1}%
\right] ^{\sigma }+\ln L_{t}+\beta \upsilon \left( b_{t+1}\right) 
\end{equation*}

\item 
\begin{eqnarray*}
c_{t}^{\sigma -1} &=&\beta v^{\prime }(b_{t+1}) \\
v^{\prime }(b_{t+1}) &=&\left( 1+r_{t+1}\right) c_{t+1}^{\sigma -1}
\end{eqnarray*}

\item We focus on \textquotedblleft balanced growth paths\textquotedblright\
where $\frac{c_{t+1}}{c_{t}}=\frac{y_{t+1}}{y_{t}}=\frac{n_{t+1}}{n_{t}}%
=\gamma $. This into the above gives 
\begin{equation*}
\gamma =\left( \beta \left( 1+r_{t}\right) \right) ^{\frac{1}{1-\sigma }}
\end{equation*}

\item Wage =%
\begin{equation*}
\left( 1-\alpha \right) AL^{-\alpha }\overset{n_{t}}{\underset{i=1}{\sum }}%
x_{it}^{\alpha }
\end{equation*}%
\begin{equation*}
w_{t}=\left( 1-\alpha \right) A\overset{n_{t}}{\underset{i=1}{\sum }}%
x_{it}^{\alpha }
\end{equation*}%
$\left( \frac{\alpha ^{2}AL^{1-\alpha }}{\phi }\right) ^{\frac{1}{1-\alpha }%
} $%
\begin{equation*}
w_{t}=\left( 1-\alpha \right) An_{t}\left( \frac{\alpha ^{2}A}{\phi }\right)
^{\frac{1}{1-\alpha }}
\end{equation*}

\item Bond market requires%
\begin{equation*}
\eta \left( 1-\varsigma \right) =\frac{1}{r}\left( \frac{\phi }{\alpha }%
-\phi \right) \left( \frac{\alpha ^{2}AL^{1-\alpha }}{\phi }\right) ^{\frac{1%
}{1-\alpha }}
\end{equation*}%
\begin{equation*}
r=\left( \frac{\phi }{\alpha }-\phi \right) \left( \frac{\alpha
^{2}AL^{1-\alpha }}{\phi }\right) ^{\frac{1}{1-\alpha }}\frac{1}{\eta \left(
1-\varsigma \right) }
\end{equation*}%
\begin{equation*}
\pi _{i}=\left( \frac{\phi }{\alpha }-\phi \right) \left( \frac{\alpha
^{2}AL^{1-\alpha }}{\phi }\right) ^{\frac{1}{1-\alpha }}
\end{equation*}

\item 
\begin{eqnarray*}
\gamma  &=&\left( \beta \left( 1+r_{t}\right) \right) ^{\frac{1}{1-\sigma }}
\\
\gamma  &=&\left( \beta \left( 1+\left( \frac{\phi }{\alpha }-\phi \right)
\left( \frac{\alpha ^{2}AL^{1-\alpha }}{\phi }\right) ^{\frac{1}{1-\alpha }}%
\frac{1}{\eta \left( 1-\varsigma \right) }\right) \right) ^{\frac{1}{%
1-\sigma }}
\end{eqnarray*}%
\begin{equation*}
\upsilon \left( b_{t}\right) =\underset{b_{t+1}}{\max }\frac{\left[
w_{t}\left( 1-l_{t}\right) \left( 1-\tau \right) +\left( 1+r_{t}\right)
b_{t}-b_{t+1}\right] ^{\sigma }}{\sigma }+\ln \left( l_{t}\right) +\beta
\upsilon \left( b_{t+1}\right) 
\end{equation*}%
\begin{equation*}
\upsilon \left( n_{t}\right) =\underset{\hat{x}_{t},n_{t+1}}{\max }\ln
\left( \overset{n_{t}}{\underset{i=1}{\sum }}\left( x_{it}^{\alpha
}-x_{it}\right) -\eta \left( n_{t+1}-n_{t}\right) \right) +\beta \upsilon
\left( n_{t+1}\right) 
\end{equation*}%
Take the FOC for $x_{it}$%
\begin{equation*}
\frac{\alpha x_{it}^{\alpha -1}-1}{c_{t}}=0
\end{equation*}%
\begin{equation*}
\alpha x_{it}^{\alpha -1}-1=0
\end{equation*}%
\begin{equation*}
\alpha x_{it}^{\alpha -1}=1
\end{equation*}%
\begin{equation*}
\alpha ^{\frac{1}{1-\alpha }}=x_{it}
\end{equation*}%
since this is the same for each good, we can write the Bellman equation as 
\begin{equation*}
\upsilon \left( n_{t}\right) =\underset{n_{t+1}}{\max }\ln \left(
n_{t}\left( \alpha ^{\frac{\alpha }{1-\alpha }}-\alpha ^{\frac{1}{1-\alpha }%
}\right) -\eta \left( n_{t+1}-n_{t}\right) \right) +\beta \upsilon \left(
n_{t+1}\right) 
\end{equation*}%
The first order condition for $n_{t+1}$ is 
\begin{equation*}
\beta \upsilon ^{^{\prime }}\left( n_{t+1}\right) =\frac{\eta }{c_{t}}
\end{equation*}%
The envelope condition for $n_{t}$ (updated) is%
\begin{equation*}
\upsilon ^{^{\prime }}\left( n_{t+1}\right) =\frac{\alpha ^{\frac{\alpha }{%
1-\alpha }}-\alpha ^{\frac{1}{1-\alpha }}+\eta }{c_{t+1}}
\end{equation*}%
Thus 
\begin{equation*}
\frac{c_{t+1}}{c_{t}}=\beta \left( \frac{\alpha ^{\frac{\alpha }{1-\alpha }%
}-\alpha ^{\frac{1}{1-\alpha }}+\eta }{\eta }\right) 
\end{equation*}%
\begin{equation*}
\gamma =\beta \left( \frac{\alpha ^{\frac{\alpha }{1-\alpha }}-\alpha ^{%
\frac{1}{1-\alpha }}+\eta }{\eta }\right) 
\end{equation*}
\end{enumerate}

********************************

Goods market clearing%
\begin{equation*}
c_{t}=\overset{n_{t}}{\underset{i=1}{\sum }}\left( x_{it}^{\alpha
}-x_{it}\right) -\eta \left( n_{t+1}-n_{t}\right) 
\end{equation*}%
bond market clearing 
\begin{equation*}
b_{t+1}-b_{t}=\eta \left( n_{t+1}-n_{t}\right) 
\end{equation*}%
\begin{equation*}
c_{t}+b_{t+1}=w_{t}\left( 1-\tau \right) +\left( 1+r_{t}\right) b_{t}
\end{equation*}%
\begin{equation*}
\overset{n_{t}}{\underset{i=1}{\sum }}\left( x_{it}^{\alpha }-x_{it}\right)
-\eta \left( n_{t+1}-n_{t}\right) +b_{t+1}=w_{t}\left( 1-\tau \right)
L+\left( 1+r_{t}\right) b_{t}
\end{equation*}%
\begin{equation*}
\overset{n_{t}}{\underset{i=1}{\sum }}\left( x_{it}^{\alpha }-x_{it}\right)
=\left( 1-\tau \right) w_{t}L+r_{t}b_{t}
\end{equation*}%
\begin{equation*}
w_{t}L=\left( 1-\alpha \right) L^{1-\alpha }\overset{n_{t}}{\underset{i=1}{%
\sum }}x_{it}^{\alpha }
\end{equation*}%
\begin{equation*}
r_{t}b_{t}=
\end{equation*}%
\begin{equation*}
L^{1-\alpha }n_{t}\left( \frac{\alpha ^{2}AL^{1-\alpha }}{\phi }\right) ^{%
\frac{\alpha }{1-\alpha }}-n_{t}\left( \frac{\alpha ^{2}AL^{1-\alpha }}{\phi 
}\right) ^{\frac{1}{1-\alpha }}=\left( 1-\alpha \right) L^{1-\alpha
}n_{t}\left( \frac{\alpha ^{2}AL^{1-\alpha }}{\phi }\right) ^{\frac{\alpha }{%
1-\alpha }}+n_{t}\left( \frac{\phi }{\alpha }-\phi \right) \left( \frac{%
\alpha ^{2}AL^{1-\alpha }}{\phi }\right) ^{\frac{1}{1-\alpha }}
\end{equation*}%
\begin{equation*}
L^{1-\alpha }n_{t}\left( \frac{\alpha ^{2}AL^{1-\alpha }}{\phi }\right) ^{%
\frac{\alpha }{1-\alpha }}-\phi n_{t}\left( \frac{\alpha ^{2}AL^{1-\alpha }}{%
\phi }\right) ^{\frac{1}{1-\alpha }}=\left( 1-\alpha \right) L^{1-\alpha
}n_{t}\left( \frac{\alpha ^{2}AL^{1-\alpha }}{\phi }\right) ^{\frac{\alpha }{%
1-\alpha }}+n_{t}\left( \frac{\phi }{\alpha }-\phi \right) \left( \frac{%
\alpha ^{2}AL^{1-\alpha }}{\phi }\right) ^{\frac{1}{1-\alpha }}
\end{equation*}%
\begin{equation*}
L^{1-\alpha }n_{t}-\phi n_{t}\frac{\alpha ^{2}AL^{1-\alpha }}{\phi }=\left(
1-\alpha \right) L^{1-\alpha }n_{t}+n_{t}\left( \frac{\phi }{\alpha }-\phi
\right) \left( \frac{\alpha ^{2}AL^{1-\alpha }}{\phi }\right) 
\end{equation*}%
\begin{equation*}
L^{1-\alpha }-\alpha ^{2}AL^{1-\alpha }=\left( 1-\alpha \right)
L_{t}^{1-\alpha }+\left( \frac{1}{\alpha }-1\right) \left( \alpha
^{2}AL^{1-\alpha }\right) 
\end{equation*}%
\begin{equation*}
1-\alpha ^{2}=1-\alpha +\left( \frac{1}{\alpha }-1\right) \alpha ^{2}
\end{equation*}
: $1-\alpha ^{2}=1-\alpha ^{2}$ : 

\end{document}
